\begin{figure*}
\centering
\resizebox{1\textwidth}{!}{
\fbox{
\begin{minipage}[c][2.2in]{\linewidth}
\begin{small}
\begin{align*}
    & x = \texttt{'The patient had presented a progressive deterioration of the general condition,}\\
    & \qquad \texttt{a fever and night sweats.'}\\
    & p = concat(t, x)\\
   & \begin{aligned}
        o = \Phi(p, \theta_h) = & \texttt{'-"fever" } \\ 
                                & \texttt{-"night sweats"'}
      \end{aligned}\\
    & \begin{aligned}
        r(o, x) = & \; [ \\ 
                         & \space\space(\texttt{The}, 0 , 3, O), (\texttt{patient}, 4 , 11, O), \dots, \\ 
                         & \space\space(\texttt{fever}, 72, 77, B_{clin}), (\texttt{and}, 79, 82, O), \\
                         & \space\space(\texttt{night}, 83, 87, B_{clin}), (\texttt{sweats}, 88, 94, I_{clin}), \dots \\
                         & ]
      \end{aligned}
\end{align*}
\end{small}
\end{minipage}}
}
    \caption{our method's prediction and structuring steps on an example. The $t$ template is illustrated in Figure \ref{fig:prompt}.}
    \label{fig:formalization}
\end{figure*}